% Part: first-order-logic
% Chapter: natural-deduction
% Section: identity

\documentclass[../../../include/open-logic-section]{subfiles}

\begin{document}

\olfileid[id]{fol}{ntd}{ide}

\olsection{\usetoken{P}{derivation} dengan \usetoken{S}{identity}}

!!^{derivation} dengan !!{identity} memerlukan aturan inferensi tambahan.

\begin{defish}
\AxiomC{}
\RightLabel{\Intro{\eq}}
\UnaryInfC{$\eq[t][t]$}
\DisplayProof
\hfill
\begin{tabular}{r}
\AxiomC{$\eq[t_1][t_2]$}
\AxiomC{$!A(t_1)$}
\RightLabel{\Elim{\eq}}
\BinaryInfC{$!A(t_2)$}
\DisplayProof
\\[3ex]
\AxiomC{$\eq[t_1][t_2]$}
\AxiomC{$!A(t_2)$}
\RightLabel{\Elim{\eq}}
\BinaryInfC{$!A(t_1)$}
\DisplayProof
\end{tabular}
\end{defish}

Dalam aturan-aturan di atas, $t$, $t_1$, dan $t_2$ merupakan suku tertutup.
Aturan \Intro{\eq} memungkinkan kita untuk !!{derive} sebarang pernyataan
identitas berbentuk $\eq[t][t]$ secara langsung, tanpa asumsi.

\begin{ex}
Jika $s$ dan $t$ merupakan suku tertutup, maka
$!A(s), \eq[s][t] \Proves !A(t)$:
\begin{prooftree}
\AxiomC{$\eq[s][t]$}
\AxiomC{$!A(s)$}
\RightLabel{$\Elim{\eq}$}
\BinaryInfC{$!A(t)$}
\end{prooftree}
Hal ini mungkin telah dikenal sebagai ``prinsip substitutabilitas hal-hal
yang identik,'' atau Hukum Leibniz.
\end{ex}

\begin{prob}
Buktikan bahwa $=$ bersifat simetris dan transitif, yakni berikan
!!{derivation} untuk $\lforall[x][\lforall[y][(\eq[x][y] \lif
    \eq[y][x])]]$ dan $\lforall[x][\lforall[y][\lforall[z][((\eq[x][y]
    \land \eq[y][z]) \lif \eq[x][z])]]]$.
\end{prob}

\begin{ex}
Kita !!{derive} !!{sentence}
\begin{align*}
& \lforall[x][\lforall[y][((!A(x) \land !A(y)) \lif \eq[x][y])]]
\intertext{dari !!{sentence}}
& \lexists[x][\lforall[y][(!A(y) \lif \eq[y][x])]]
\end{align*}
Misalkan tiga konstanta yang digunakan di bawah ini adalah simbol-simbol
baru yang berbeda dan tidak muncul dalam formula matriks (dan dengan
demikian tidak muncul dalam premis).
Kita mengembangkan !!{derivation} tersebut secara mundur:
\begin{prooftree}
\AxiomC{$\lexists[x][\lforall[y][(!A(y) \lif \eq[y][x])]]
\quad \Discharge{!A(a) \land !A(b)}{1}$}
\DeduceC{$\eq[a][b]$}
\DischargeRule{\Intro{\lif}}{1}
\UnaryInfC{$((!A(a) \land !A(b)) \lif \eq[a][b])$}
\RightLabel{\Intro{\lforall}}
\UnaryInfC{$\lforall[y][((!A(a) \land !A(y)) \lif \eq[a][y])]$}
\RightLabel{\Intro{\lforall}}
\UnaryInfC{$\lforall[x][\lforall[y][((!A(x) \land !A(y)) \lif \eq[x][y])]]$}
\end{prooftree}
Sekarang kita harus menggunakan asumsi utama: karena asumsi itu merupakan
!!{formula} eksistensial, kita menggunakan \Elim{\lexists} untuk
!!{derive} kesimpulan antara $\eq[a][b]$.
\begin{prooftree}
\AxiomC{$\lexists[x][\lforall[y][(!A(y) \lif \eq[y][x])]]$}
\AxiomC{$\Discharge{\lforall[y][(!A(y) \lif \eq[y][c]])}{2}$}
\noLine
\UnaryInfC{$\Discharge{!A(a) \land !A(b)}{1}$}
\DeduceC{$\eq[a][b]$}
\DischargeRule{\Elim{\lexists}}{2}
\BinaryInfC{$\eq[a][b]$}
\DischargeRule{\Intro{\lif}}{1}
\UnaryInfC{$((!A(a) \land !A(b)) \lif \eq[a][b])$}
\RightLabel{\Intro{\lforall}}
\UnaryInfC{$\lforall[y][((!A(a) \land !A(y)) \lif \eq[a][y])]$}
\RightLabel{\Intro{\lforall}}
\UnaryInfC{$\lforall[x][\lforall[y][((!A(x) \land !A(y)) \lif \eq[x][y])]]$}
\end{prooftree}
Sub-!!{derivation} di kanan atas diselesaikan dengan menggunakan
asumsi-asumsinya untuk menunjukkan $\eq[a][c]$ dan $\eq[b][c]$. Hal ini
memerlukan dua !!{derivation} terpisah. !!^{derivation} untuk
$\eq[a][c]$ adalah sebagai berikut:
\begin{prooftree}
\AxiomC{$\Discharge{\lforall[y][(!A(y) \lif \eq[y][c]])}{2}$}
\RightLabel{\Elim{\lforall}}
\UnaryInfC{$!A(a) \lif \eq[a][c]$}
\AxiomC{$\Discharge{!A(a) \land !A(b)}{1}$}
\RightLabel{\Elim{\land}}
\UnaryInfC{$!A(a)$}
\RightLabel{\Elim{\lif}}
\BinaryInfC{$\eq[a][c]$}
\end{prooftree}
Dari $\eq[a][c]$ dan $\eq[b][c]$ kita !!{derive} $\eq[a][b]$ dengan
\Elim{\eq}.
\end{ex}

\begin{prob}
Berikan !!{derivation} untuk !!{formula} berikut:
\begin{enumerate}
\item $\lforall[x][\lforall[y][((\eq[x][y] \land !A(x)) \lif !A(y))]]$
\item $\lexists[x][!A(x)] \land \lforall[y][\lforall[z][((!A(y) \land
    !A(z)) \lif \eq[y][z])]] \lif \lexists[x][(!A(x) \land
  \lforall[y][(!A(y) \lif \eq[y][x])])]$
\end{enumerate}
\end{prob}

\end{document}
